% Blatt 7, A3. Aufgabenstellung: aufgaben/Blatt07.pdf, S. 2.

\subsection*{Gegeben und Modell}
Die bis zum Rand gefüllte Mauer ist \(H=\qty{50}{\metre}\) hoch. Die
Seeoberfläche wird von \(\qty{0.80}{\kilo\metre\squared}\) auf
\(\qty{1.2}{\kilo\metre\squared}\) vergrössert; der Wasserstand bleibt
gleich. Gesucht sind die örtlichen Drücke am Fuss und auf halber Mauerhöhe.

Wir nehmen ruhendes Wasser mit konstanter Dichte
\(\rho=\qty{1000}{\kilogram\per\metre\cubed}\) und
\(g=\qty{9.81}{\metre\per\second\squared}\) an. Diese Näherungswerte
sind nicht auf dem Aufgabenblatt vorgegeben. Der Luftdruck an der freien
Oberfläche sei \(p_0\). Mit der Tiefe \(h\) gilt \(p=p_0+\rho gh\).

\subsection*{Druck am Fuss und auf halber Höhe}
Der Fuss liegt \(H\), die halbe Höhe \(H/2\) unter dem Wasserspiegel. Für
den Schweredruck, also den Überdruck gegenüber \(p_0\), ergibt sich
\begin{align}
  \Delta p_{\mathrm{Fuss}}&=\rho gH
    =\qty{490500}{\pascal}\approx\qty{491}{\kilo\pascal},\\
  \Delta p_{\mathrm{halb}}&=\rho g\frac{H}{2}
    =\qty{245250}{\pascal}\approx\qty{245}{\kilo\pascal}.
\end{align}
Die zugehörigen absoluten Drücke sind \(p_0+\qty{490500}{\pascal}\) und
\(p_0+\qty{245250}{\pascal}\). Wählt man \emph{beispielhaft} den
Normluftdruck \(p_0=\qty{101325}{\pascal}\), so ergeben sich etwa
\(\qty{592}{\kilo\pascal}\) beziehungsweise \(\qty{347}{\kilo\pascal}\) absolut.
Der tatsächliche Luftdruck am Stausee ist nicht angegeben.

Steht auf der trockenen Mauerseite derselbe Luftdruck an, sind die obigen
Überdrücke zugleich die lokalen Druckdifferenzen über die Mauer. Andernfalls
ist jeweils der tatsächliche äussere Druck vom inneren absoluten Druck
abzuziehen.

\subsection*{Vergrösserung des Sees}
\begin{ergebnis}[Unveränderter Druck]
An beiden Stellen bleiben die Schweredrücke und bei unverändertem
Oberflächen- und Aussendruck auch die absoluten Drücke und Druckdifferenzen
unverändert. Die Vergrösserung der Seeoberfläche ändert die jeweilige Tiefe
nicht.
\end{ergebnis}
Die zusätzlich enthaltene Wassermenge bestimmt nicht den örtlichen Druck bei
vorgegebenem Füllstand. Daraus folgt allein keine Aussage über eine veränderte
Gesamtkraft, falls sich etwa die benetzte Mauerfläche ändern sollte.

\begin{kontrolle}[Plausibilität]
Die halbe Tiefe ergibt den halben Schweredruck. Die Einheit
\(\rho gh\) ist \(\si{\newton\per\metre\squared}\). Bei \(h=0\) ist
der Schweredruck null, der absolute Druck aber \(p_0\).
\end{kontrolle}

\begin{merksatz}
Entscheidend für den hydrostatischen Druck sind die vertikale Tiefe und der
Druck an der Oberfläche, nicht die gesamte Wassermenge.
\end{merksatz}
